\subsection*{Passes 4753--4760: residue algebra, odd-$q$ distance, dependency cubes, and the 24-dimensional boundary}
\addcontentsline{toc}{subsection}{Passes 4753--4760: algebra, odd-q distance, dependency cubes, and 24D boundary}

\paragraph{Pass 4753: the complete residue orbital algebra.}
For the transitive $270$-residue action of $G=\PSp(4,3)$, the orbital algebra
$\mathcal A_{\mathbb Q}$ has dimension $12$ and center dimension $9$.  Exact
structure constants and a split rank-$20$ idempotent in the unique
noncommutative block give
\[
 \boxed{\mathcal A_{\mathbb Q}\cong
 \mathbb Q^6\times\mathbb Q(\sqrt{-3})\times M_2(\mathbb Q)},
 \qquad
 \mathcal A_{\mathbb C}\cong\mathbb C^8\times M_2(\mathbb C).
\]
The $M_2$ block is the multiplicity-two degree-$20$ constituent already
responsible for the cold-router eigenvalues $1\pm\sqrt{13}$.  A primitive
rank-$24$ central projector is also obtained exactly.

\paragraph{Pass 4754: the odd-$q$ line-kernel distance theorem.}
Let $A_*(q)$ be the line-intersection adjacency matrix of $W(3,q)$.  For every
odd prime power $q$, a generalized-quadrangle double count gives
$d(\ker_{\mathbb F_2}A_*)\ge q+1$.  In the dual parabolic quadrangle
$Q(4,q)$, an anisotropic-plane conic supplies a kernel word of weight $q+1$.
Hence
\[
 \boxed{d(\ker_{\mathbb F_2}A_*(q))=q+1\quad(q\ \text{odd}).}
\]
The canonical projective $J$-involution fixes $q+1$ lines when
$q\equiv3\pmod4$, but $q+3$ when $q\equiv1\pmod4$.  Thus the distance law is
uniform while the involution/minimum-shell coincidence has a quadratic-character
branch.  Related odd-$q$ binary codes appear in the finite-geometry literature;
no publication-novelty claim is attached without a full prior-art comparison.

\paragraph{Pass 4755: representation bridge versus Leech-lattice firewall.}
The rank-$24$ projector sends the $270$ residues to an integral Gram orbit after
normalization to norm $24$, with pair inner products $9,4,-1,-6,-16$.  The
lattice generated by the orbit has
\[
 \boxed{\det L=2^{10}3^{10}5^{23}}.
\]
Because the Gram contains an entry $1$, every scalar rescaling that remains
integral has integral scalar and therefore preserves the odd parity of the
$5$-adic determinant valuation.  Consequently this \emph{canonical projector
lattice} cannot be a full-rank sublattice of an integral unimodular
$24$-lattice such as Leech.  This does not exclude another commensurable
$G$-invariant lattice in the same rational representation.

\paragraph{Pass 4756: higher dependency geometry.}
The cold graph has $1080$ triangles: $540$ XOR-zero dependency circuits and
$540$ nondependency triangles.  Every cold edge belongs to exactly one triangle
of each species.  The dependency code $[270,240,3]$ begins
\[
 A_3=540,\qquad A_4=3645=1620+1620+405,\qquad A_5=29376,
\]
where the displayed decomposition is the exact $G$-orbit split of the weight-$4$
shell.

\paragraph{Pass 4758: the triangle incidence is $135Q_3$ and reconstructs selected135.}
Join a dependency triangle to a nondependency triangle when they share a cold
edge.  The resulting cubic bipartite graph is
\[
 \boxed{135\,Q_3}.
\]
Each cube uses six residues forming $K_6-3K_2$; their W33 support union is an
eight-line word of $H_{10}^{\perp}$.  The $135$ such words form one
$G$-orbit.  Their support-intersection-$4$ graph is explicitly
$G$-equivariantly identical to the selected135 graph, and the six residues in
each cube map to exactly the six selected270 lines incident with the
corresponding selected135 vertex.  Thus the dependency geometry reconstructs
the full $135\leftrightarrow270$ incidence layer.

\paragraph{Pass 4759: a $45$-vector spherical quotient of the same $24$-space.}
Sum the six rank-$24$ residue vectors in each cube.  The $135$ sums collapse in
three-element packets to $45$ distinct vectors.  After dividing their Gram by
$9$, they have norm $16$ and mutual inner products $-4$ or $1$.  If $A_{45}$
is the $-4$ relation, then
\[
 \boxed{G_{45}=15I-5A_{45}+J},
 \qquad A_{45}\text{ is }\operatorname{SRG}(45,12,3,3).
\]
The Gram has rank $24$ and nonzero eigenvalue $30$, giving the canonical
$24$-dimensional spherical embedding of the $45$-point quotient directly from
the residue dependencies.

\paragraph{Pass 4760: the binary $24$-dimensional intersection is now explicit.}
Let $C_\star$ be the binary code on the forty W33 line coordinates generated by
the forty point-stars.  Then $\dim C_\star=25$, while
$H_{10}^{\perp}$ has dimension $30$ and the combined rank is $31$.  The $135$
cube-union words lie in both codes and span rank $24$, so
\[
 \boxed{\operatorname{span}\{\text{cube unions}\}
 =C_\star\cap H_{10}^{\perp}=[40,24,6].}
\]
Its complete enumerator begins
$1+240z^6+1485z^8+27792z^{10}+169600z^{12}+\cdots+z^{40}$.
The $240$ minimum words are exactly the sums of point-stars at the $240$
collinear W33 point pairs.  No cross-characteristic identification with the
rank-$24$ characteristic-zero constituent of Pass~4753 is inferred from the
shared dimension.

\paragraph{Pass 4757: exact radius-six parity-Hamiltonian sectors.}
For $\mathcal H=-\sum_R\prod_{i\in R}Z_i$, the syndrome/coset-leader census is
exhaustive through original error weight six.  At weight six there are
$3{,}819{,}900$ distinct syndromes and exactly $18{,}480$ doubleton classes,
with no triples.  Two skew line defects cost $96$, giving binding $-12$
relative to two isolated $54$-cost defects; meeting defects cost $108$.  For
independent line flips of probability $p$, writing $a=1-2p$ gives
\[
 \mathbb E[\mathcal H]=-270a^4,
\]
\[
 \operatorname{Var}(\mathcal H)=270(1-a^8)+21600(a^6-a^8)+3240(a^4-a^8).
\]
This last statement is an exact classical independent-noise calculation, not a
thermal-equilibrium or microscopic-physics claim.

\begin{quote}\small
\textbf{Evidence boundary.}  All finite claims above are accompanied by executable
reconstruction scripts.  The Leech statement is a no-go for the canonical
projected lattice only; the all-odd-$q$ result is a finite-geometry theorem with
an explicit prior-art firewall; and the Hamiltonian statements remain within
the commuting parity-check model.
\end{quote}
